% !TEX root = ../main.tex
Modern latency-critical RPC services often run on commodity multi-core
servers with one worker pinned to each CPU core. To avoid centralized
queue contention and preserve cache locality, each worker maintains a
local request queue, and incoming packets are distributed by RSS, flow
hashing, or a user-space ingress thread. This design scales well under
balanced load, but it can create transient per-core queue imbalance
under bursty arrivals, skewed flow distribution, and heterogeneous
service times. For microsecond-scale RPC, even a short period of local
queue buildup can cause P99 or P99.9 tail-latency SLO violations.

Existing work stealing improves work conservation by letting idle cores
pull tasks from busy queues; work shedding and threshold-based migration
react when a queue becomes overloaded. However, these mechanisms mainly
answer one question: \emph{which core has too much work?} They do not
explicitly answer three finer-grained SLO questions:
(i) which queued request will miss its SLO if it stays local;
(ii) which of those requests can still meet its SLO after migration;
(iii) which target core can accept the migrated request without becoming
risky itself.

\note{Sharpen the gap statement and add one concrete motivating number
from the evaluation here.}

\textbf{Contributions.} This paper makes three contributions:
\begin{itemize}
  \item We define \emph{task-level rescuability} for microsecond-scale
  RPC scheduling, using local-vs-remote completion-time counterfactuals
  under a migration-cost model (\S\ref{sec:model}).
  \item We design \emph{migration-safe target selection}, which prevents
  cross-core migration from transferring SLO risk from the source core
  to the target core (\S\ref{sec:design}).
  \item We implement RescueSched as a commodity user-space prototype and
  show that selective descriptor migration reduces SLO violations and
  useless migrations compared with work stealing, work shedding, and
  threshold-based migration baselines (\S\ref{sec:eval}).
\end{itemize}